
% !TeX program = pdflatex
% !TeX program = lualatex
% !BIB program = biber
% !TeX encoding = UTF-8
% !TeX spellcheck = en_US
% !TeX root = BookChap0.tex
% !TeX TS-program = pdflatex

\documentclass[
    11pt,
    fleqn,
    a4paper
]{LegrandOrangeBook}
%%%%%%%%%%%%%%%%%%%%%%%%%%%%%%%%%%%%%%%%%%%%%%%%%%%%%%%%%%%%%%%%

\usepackage{url}  % Solo para \url si lo necesitas

\title{\Huge\textbf{A General Equation of Science}\\\Large\textbf{Implications}}

\usepackage{float}
%\usepackage{amsmath,amssymb}
\usepackage{tcolorbox}
\definecolor{ocre}{RGB}{243, 102, 25} % Color for highlighting
\chapterimage{orange1.jpg} % Chapter heading image
\chapterspaceabove{6.5cm} % Whitespace to chapter title
\chapterspacebelow{6.75cm} % Whitespace to text on chapter pages


\usepackage{fontspec}
\setmainfont{Times New Roman}   % puedes cambiar a Calibri, Georgia, Libertinus Serif,

\usepackage{fancyvrb}

% ---------------------------------------------------------------------------------------
% Other mathematical packages
\usepackage{bm}
\usepackage{centernot}
\usepackage{xkeyval}
\usepackage{placeins}
\usepackage{listings}
% Algorithm packages
\usepackage{algorithm}
\usepackage{algpseudocode}
\newtheorem{assumption}{Assumption}

\usepackage{tocloft}

\setlength{\cftfignumwidth}{4em}
\setlength{\cftfigindent}{1em}

\setlength{\cftpartnumwidth}{3em}

\usepackage{minted}
\usemintedstyle{default} % o vs, tango, friendly, etc.

\usepackage{pgfplots}
\usepackage{tikz}
\usetikzlibrary{fit}
\usetikzlibrary{shapes,arrows,positioning,calc}
\usetikzlibrary{patterns, decorations.pathmorphing, backgrounds, shapes, arrows, positioning, calc, 3d}

\usepackage{comment}
\usepackage{caption}
\usepackage{subcaption}

% BIBLIOGRAPHY CONFIGURATION - SIMPLIFIED
% Remove one of the duplicate lines and use only this:
\addbibresource{bibliography.bib}

% ===== MDFRAMED ENVIRONMENTS (SIMPLIFIED) =====
\usepackage{mdframed}
\usepackage{xcolor}
%\definecolor{brown}{RGB}{165,42,42}
%%%%%%%%%%%%%%%%%%%%%%%%%%%%%%%%%%%%%%%%%%
% AÑADIDO SOLO PARA complements.tex
%%%%%%%%%%%%%%%%%%%%%%%%%%%%%%%%%%%%%%%%%%
\usepackage{amsmath, amssymb}
%\usepackage{amsthm}
\usepackage{mathrsfs}
\usepackage{mathtools}
\usepackage{booktabs}


%\begin{comment}
%\newtheorem{definition}{Definition}
%\newtheorem{theorem}{Theorem}
%\newtheorem{corollary}{Corollary}
%\newtheorem{remark}{Remark}
%\newtheorem{example}{Example}
%\newtheorem{proof}{Proof}
%\end{comment}

%\newtheorem{lemma}{Lemma}
\newtheorem{application}{Application}
\newtheorem{myremark}{Remark}[chapter]

\begin{comment}
\newcommand{\R}{\mathbb{R}}
\newcommand{\N}{\mathcal{N}}
\end{comment}


\newcommand{\minitocboxed}{%
  \begingroup
  \tcbstart
  \begin{tcolorbox}[colback=orange!8, arc=3pt, boxrule=1pt]
    \minitoc
  \end{tcolorbox}
  \tcbend
  \endgroup
}

\newcommand{\E}{\mathbb{E}}
\newcommand{\Var}{\operatorname{Var}}
\newcommand{\Cov}{\operatorname{Cov}}
\newcommand{\Prob}{\mathbb{P}}
\newcommand{\as}{\text{ a.s.}}
\newcommand{\iid}{\stackrel{\text{i.i.d.}}{\sim}}
\newcommand{\convp}{\xrightarrow{p}}
\newcommand{\convas}{\xrightarrow{\text{a.s.}}}
\newcommand{\convd}{\xrightarrow{d}}
\newcommand{\plim}{\operatorname*{plim}}
\newcommand{\ind}{\mathbf{1}}
%%%%%%%%%%%%%%%%%%%%%%%%%%%%%%%%%%%%%%%%%%

% Problem box
\newmdenv[
    skipabove=7pt,
    skipbelow=7pt,
    backgroundcolor=brown!30,
    linecolor=ocre!50,
    linewidth=1pt,
    roundcorner=3pt,
    innerleftmargin=8pt,
    innerrightmargin=8pt,
    innertopmargin=8pt,
    innerbottommargin=8pt
]{problembox}

% Description box
\newmdenv[
    skipabove=7pt,
    skipbelow=7pt,
    rightline=false,
    leftline=true,
    topline=false,
    bottomline=false,
    linecolor=ocre!50,
    linewidth=4pt,
    backgroundcolor=orange!30,
    innerleftmargin=5pt,
    innerrightmargin=5pt,
    innertopmargin=5pt,
    innerbottommargin=5pt
]{descriptionbox}

% Shadow box
\newmdenv[
    skipabove=7pt,
    skipbelow=7pt,
    backgroundcolor=brown!8,
    linecolor=ocre!50,
    linewidth=1pt,
    innerleftmargin=8pt,
    innerrightmargin=8pt,
    innertopmargin=8pt,
    innerbottommargin=8pt
]{shadowbox}

% Code box
\newmdenv[
    skipabove=7pt,
    skipbelow=7pt,
    backgroundcolor=brown!8,
    linecolor=ocre!50,
    linewidth=1pt,
    innerleftmargin=8pt,
    innerrightmargin=8pt,
    innertopmargin=8pt,
    innerbottommargin=8pt,
    font=\small
]{codeboxbase}

\newenvironment{codebox}
{\VerbatimEnvironment
 \begin{codeboxbase}%
 \begin{Verbatim}}
{\end{Verbatim}%
 \end{codeboxbase}}

% Example box
\newmdenv[
    skipabove=7pt,
    skipbelow=7pt,
    backgroundcolor=orange!8,
    linecolor=ocre!50,
    linewidth=1pt,
    innerleftmargin=8pt,
    innerrightmargin=8pt,
    innertopmargin=8pt,
    innerbottommargin=8pt
]{examplebox}

% Remark box
\newenvironment{remarkbox}
{\begin{mdframed}[
   backgroundcolor=blue!10,
   linecolor=blue,
   linewidth=1pt]
   \textbf{Remark: }\ignorespaces
}
{\end{mdframed}}



% Algorithm box - simplified
\newmdenv[
    skipabove=7pt,
    skipbelow=7pt,
    backgroundcolor=yellow!8,
    linecolor=ocre!50,
    linewidth=1pt,
    innerleftmargin=8pt,
    innerrightmargin=8pt,
    innertopmargin=8pt,
    innerbottommargin=8pt
]{algorithmbox}

\newenvironment{importantbox}[1][]
  {\begin{mdframed}[
     backgroundcolor=orange!10,
     linecolor=orange,
     linewidth=1pt,
     frametitle={#1}]
  }
  {\end{mdframed}}

% Custom mathematical commands
\newcommand{\R}{\mathbb{R}}
\newcommand{\N}{\mathbb{N}}
\newcommand{\Z}{\mathbb{Z}}
\newcommand{\Q}{\mathbb{Q}}
\newcommand{\C}{\mathbb{C}}

\DeclareMathOperator{\rank}{rank}

% METADATA FOR PDF - MOVED HERE
\hypersetup{
    pdftitle={A General Equation of Science. Implications},
    pdfauthor={Enrique Castillo Ron and Alfonso Fernández-Canteli},
    pdfsubject={General Equation of Science},
    pdfkeywords={Science, Equation, Modeling, Foundations},
}

\usepackage[tight]{minitoc}
\dominitoc
\nomtcrule  % Sin líneas horizontales
\setlength{\mtcindent}{0pt}  % Sin sangría
\renewcommand{\mtcfont}{\small\sffamily}  % Fuente pequeña sans-serif
\renewcommand{\mtcSfont}{\small\sffamily}  % Para secciones

% Entorno shadowbox personalizado
\newenvironment{pythoncode}
{\begin{shadowbox}\begin{lstlisting}[language=Python]}
{\end{lstlisting}\end{shadowbox}}


%\includeonly{cap04a_MethodsSolvingFE}


\begin{document}
\setcounter{minitocdepth}{1}


%----------------------------------------------------------------------------------------
%	TITLE PAGE
%----------------------------------------------------------------------------------------
\titlepage
    {\includegraphics[width=210mm,height=297mm]{background.png}}
    {
        \centering\sffamily
        {\Huge\bfseries A General Equation of Science\par}
        \vspace{16pt}
        {\LARGE Implications\par}
        \vspace{24pt}
        {\huge\bfseries Enrique Castillo Ron\\ and\\ Alfonso Fernández-Canteli\par}
    }

%----------------------------------------------------------------------------------------
%	COPYRIGHT PAGE
%----------------------------------------------------------------------------------------

\thispagestyle{empty}

~\vfill

\noindent Copyright \copyright\ 2025 Enrique Castillo Ron and Alfonso Fernández-Canteli\\

\noindent \textsc{Published by Publisher}\\

\noindent \textsc{\href{https://www.latextemplates.com/template/legrand-orange-book}{book-website.com}}\\

\noindent Licensed under the Creative Commons Attribution-NonCommercial 4.0 License (the ``License''). You may not use this file except in compliance with the License. You may obtain a copy of the License at \url{https://creativecommons.org/licenses/by-nc-sa/4.0}. Unless required by applicable law or agreed to in writing, software distributed under the License is distributed on an \textsc{``as is'' basis, without warranties or conditions of any kind}, either express or implied. See the License for the specific language governing permissions and limitations under the License.\\

\noindent \textit{First printing, March 2025}



%----------------------------------------------------------------------------------------
%	TABLE OF CONTENTS
%----------------------------------------------------------------------------------------

\pagestyle{empty}
%\listoffigures % Comment if not needed
%\listoftables % Comment if not needed

\pagestyle{fancy}

\tableofcontents
\adjustmtc
\cleardoublepage

\begin{comment}
\renewenvironment{proof}{%
  \begin{pBox}%
  \begin{proofT}%
}{%
  \end{proofT}%
  \end{pBox}%
}
\end{comment}
\setcounter{mtc}{0}
%\setcounter{minitocdepth}{1}
%\faketableofcontents  % Reinicia completamente minitoc
%----------------------------------------------------------------------------------------
%	PART I: WHAT IS THIS BOOK ABOUT?
%----------------------------------------------------------------------------------------

\chapterimage{orange2.jpg}
\chapterspaceabove{6.75cm}
\chapterspacebelow{7.25cm}

\part{What is This Book About}
\parttoc



%\begin{comment}
\input{cap01_intro}
\input{cap02_historia}
\input{cap03_buckingham}
\input{cap04_funcionales}
%\input{cap04a_MethodsSolvingFEOLD}
\input{cap04a_MethodsSolvingFE}
%\input{cap04a_MethodsSolvingFEClaude}
%\input{cap04a_MethodsSolvingFEchatGPT}
%\input{cap04a_MethodsSolvingFEGrok}
%\input{cap04a_MethodsSolvingFEDeepseek}


\input{cap05_TranDimensionlessForm}




\input{cap05a_FormBuckFunctions}
%\input{cap05b_forma_buckingham}




\input{cap06_copulas_intro}
\input{cap07_copulas_arquimedianas}




%\begin{comment}
\input{cap08_redes_bayesianas_v2}
\input{cap09_inferencia_bayesiana}
\input{cap10_fatiga}
%\end{comment}
\input{cap10a_trivariateWeibullHyperbolaModel}

\input{cap11_acciones_interacciones}



\input{cap12_redes_funcionales}
\input{cap13_herramientas}



\input{cap14_conclusiones}
\input{cap15_futuro}



\appendix
\input{apendiceA_probabilidad}



\input{apendiceB_codigo}
\input{apendiceC_demostraciones}
\printbibliography[title={Bibliography}]
\end{document} 